% GENERATED FILE. Edit canonical lessons, then run npm run generate:books.
% canonical-source-hash: fnv1a64:d39e4d0eb362a9e0
% canonical-lessons: HI-C393-saphalta, HI-C393-upalabdhi, HI-C393-hasilkarna, HI-C393-taraqqi, HI-C393-gaurav

\chapter{Success}
\label{ch:hi-success}
\hlchaptermodality{393}

\begin{chapteropening}
\textbf{By the end of this chapter:} \emph{I can say what it will take to succeed.}

Complete the last lesson of chapter 393: I can say what it will take to succeed.
\end{chapteropening}

\section[\texorpdfstring{\hi{सफलता} (saphaltā)}{सफलता (saphaltā)}]{\hi{सफलता} (saphaltā) — success}

\label{lesson:HI-C393-saphalta}



\begin{quote}
Before the new one: say the Hindi for experience, then the Hindi for to occur to.
\end{quote}

\subsection*{You'll want to know: \hi{सफलता}}
\textbf{\hi{सफलता}} — \emph{saphaltā} — \textquotedblleft{}success\textquotedblright{}.

\begin{quote}
\hi{अगर} \hi{तुम} \hi{हार} \hi{नहीं} \hi{मानोगे} \hi{तो} \hi{सफलता} \hi{ज़रूर} \hi{मिलेगी।}
\end{quote}

\begin{quote}
\emph{agar tum hār nahīṃ mānoge to saphaltā zarūr milegī.}
\end{quote}

\begin{quote}
If you don't give up, success will surely come.
\end{quote}

Feminine.

\begin{cousinweb}
\textbf{\hi{सफलता}} (\emph{saphaltā}) is Sanskrit \emph{saphala}, \textquotedblleft{}bearing fruit; having good results\textquotedblright{} (Monier-Williams), with the abstract ending \emph{-tā}.
\end{cousinweb}

\begin{grammarlens}[title={What you've built}]
One more word for saying what happens if.
\end{grammarlens}

\subsection*{Your turn}
\begin{itemize}
\raggedright
  \item \emph{Say it:} \emph{saphaltā}
  \item \emph{Say it:} \emph{saphaltā}, once more
  \item \emph{Say it:} \emph{agar tum hār nahīṃ mānoge to saphaltā zarūr milegī.} — If you don't give up, success will surely come.
  \item \emph{Recall:} say the Hindi for experience, then the Hindi for to occur to, then say \emph{saphaltā} again
\end{itemize}

\begin{tcolorbox}[breakable,colback=teal!4,colframe=teal!35!black,title={Before you move on}]
What does \hi{सफलता} mean? (\textquotedblleft{}Success\textquotedblright{}.) Say it once more.
\end{tcolorbox}
\section[\texorpdfstring{\hi{उपलब्धि} (upalabdhi)}{उपलब्धि (upalabdhi)}]{\hi{उपलब्धि} (upalabdhi) — an achievement}

\label{lesson:HI-C393-upalabdhi}



\begin{quote}
Before the new one: say the Hindi for to occur to, then the Hindi for success.
\end{quote}

\subsection*{You'll want to know: \hi{उपलब्धि}}
\textbf{\hi{उपलब्धि}} — \emph{upalabdhi} — \textquotedblleft{}an achievement\textquotedblright{}.

\begin{quote}
\hi{अगर} \hi{गाँव} \hi{में} \hi{बिजली} \hi{आई} \hi{तो} \hi{यह} \hi{हमारी} \hi{सबसे} \hi{बड़ी} \hi{उपलब्धि} \hi{होगी।}
\end{quote}

\begin{quote}
\emph{agar gā̃v meṁ bijlī āī to yah hamārī sabse baṛī upalabdhi hogī.}
\end{quote}

\begin{quote}
If electricity reaches the village, it will be our greatest achievement.
\end{quote}

Feminine.

\begin{cousinweb}
\textbf{\hi{उपलब्धि}} (\emph{upalabdhi}) is a tatsama, borrowed from Sanskrit \emph{upalabdhi}, \textquotedblleft{}obtaining, perceiving\textquotedblright{} (Monier-Williams).
\end{cousinweb}

\begin{grammarlens}[title={What you've built}]
One more word for saying what happens if.
\end{grammarlens}

\subsection*{Your turn}
\begin{itemize}
\raggedright
  \item \emph{Say it:} \emph{upalabdhi}
  \item \emph{Say it:} \emph{upalabdhi}, once more
  \item \emph{Say it:} \emph{agar gā̃v meṁ bijlī āī to yah hamārī sabse baṛī upalabdhi hogī.} — If electricity reaches the village, it will be our greatest achievement.
  \item \emph{Recall:} say the Hindi for to occur to, then the Hindi for success, then say \emph{upalabdhi} again
\end{itemize}

\begin{tcolorbox}[breakable,colback=teal!4,colframe=teal!35!black,title={Before you move on}]
What does \hi{उपलब्धि} mean? (\textquotedblleft{}An achievement\textquotedblright{}.) Say it once more.
\end{tcolorbox}
\section[\texorpdfstring{\hi{हासिल} \hi{करना} (hāsil karnā)}{हासिल करना (hāsil karnā)}]{\hi{हासिल} \hi{करना} (hāsil karnā) — to attain, to achieve}

\label{lesson:HI-C393-hasilkarna}



\begin{quote}
Before the new one: say the Hindi for success, then the Hindi for an achievement.
\end{quote}

\subsection*{You'll want to know: \hi{हासिल} \hi{करना}}
\textbf{\hi{हासिल} \hi{करना}} — \emph{hāsil karnā} — \textquotedblleft{}to attain, to achieve\textquotedblright{}.

\begin{quote}
\hi{अगर} \hi{तुम} \hi{रोज़} \hi{अभ्यास} \hi{करोगे} \hi{तो} \hi{अपना} \hi{लक्ष्य} \hi{हासिल} \hi{करोगे।}
\end{quote}

\begin{quote}
\emph{agar tum roz abhyās karoge to apnā lakṣya hāsil karoge.}
\end{quote}

\begin{quote}
If you practise every day, you'll reach your goal.
\end{quote}

Past, with \hi{ने} (he or she …): \hi{उसने} \hi{अपना} \hi{लक्ष्य} \hi{हासिल} \hi{किया}, \emph{usne apnā lakṣya hāsil kiyā}.

\begin{cousinweb}
\textbf{\hi{हासिल} \hi{करना}} (\emph{hāsil karnā}) joins \hi{हासिल} to \hi{करना}, \textquotedblleft{}to do\textquotedblright{}. \hi{हासिल} is usually traced to Arabic \emph{ḥāṣil}, \textquotedblleft{}resulting, obtained\textquotedblright{}, which reached Hindi through Persian.
\end{cousinweb}

\begin{grammarlens}[title={What you've built}]
One more word for saying what happens if.
\end{grammarlens}

\subsection*{Your turn}
\begin{itemize}
\raggedright
  \item \emph{Say it:} \emph{hāsil karnā}
  \item \emph{Say it:} \emph{hāsil karnā}, once more
  \item \emph{Say it:} \emph{agar tum roz abhyās karoge to apnā lakṣya hāsil karoge.} — If you practise every day, you'll reach your goal.
  \item \emph{Recall:} say the Hindi for success, then the Hindi for an achievement, then say \emph{hāsil karnā} again
\end{itemize}

\begin{tcolorbox}[breakable,colback=teal!4,colframe=teal!35!black,title={Before you move on}]
What does \hi{हासिल} \hi{करना} mean? (\textquotedblleft{}To attain, to achieve\textquotedblright{}.) Say it once more.
\end{tcolorbox}
\section[\texorpdfstring{\hi{तरक़्क़ी} (taraqqī)}{तरक़्क़ी (taraqqī)}]{\hi{तरक़्क़ी} (taraqqī) — progress; a promotion}

\label{lesson:HI-C393-taraqqi}



\begin{quote}
Before the new one: say the Hindi for an achievement, then the Hindi for to attain.
\end{quote}

\subsection*{You'll want to know: \hi{तरक़्क़ी}}
\textbf{\hi{तरक़्क़ी}} — \emph{taraqqī} — \textquotedblleft{}progress; a promotion\textquotedblright{}.

\begin{quote}
\hi{अगर} \hi{तुम} \hi{अच्छा} \hi{काम} \hi{करोगे} \hi{तो} \hi{तुम्हें} \hi{तरक़्क़ी} \hi{मिलेगी।}
\end{quote}

\begin{quote}
\emph{agar tum acchā kām karoge to tumhẽ taraqqī milegī.}
\end{quote}

\begin{quote}
If you do good work, you'll get a promotion.
\end{quote}

Feminine.

\begin{cousinweb}
\textbf{\hi{तरक़्क़ी}} (\emph{taraqqī}) is usually traced to Arabic \emph{taraqqī}, \textquotedblleft{}rising, ascent\textquotedblright{}, which reached Hindi through Persian.
\end{cousinweb}

\begin{grammarlens}[title={What you've built}]
One more word for saying what happens if.
\end{grammarlens}

\subsection*{Your turn}
\begin{itemize}
\raggedright
  \item \emph{Say it:} \emph{taraqqī}
  \item \emph{Say it:} \emph{taraqqī}, once more
  \item \emph{Say it:} \emph{agar tum acchā kām karoge to tumhẽ taraqqī milegī.} — If you do good work, you'll get a promotion.
  \item \emph{Recall:} say the Hindi for an achievement, then the Hindi for to attain, then say \emph{taraqqī} again
\end{itemize}

\begin{tcolorbox}[breakable,colback=teal!4,colframe=teal!35!black,title={Before you move on}]
What does \hi{तरक़्क़ी} mean? (\textquotedblleft{}Progress; a promotion\textquotedblright{}.) Say it once more.
\end{tcolorbox}
\section[\texorpdfstring{\hi{गौरव} (gaurav)}{गौरव (gaurav)}]{\hi{गौरव} (gaurav) — pride, a matter of honour}

\label{lesson:HI-C393-gaurav}



\begin{quote}
Before the new one: say the Hindi for to attain, then the Hindi for a promotion.
\end{quote}

\subsection*{You'll want to know: \hi{गौरव}}
\textbf{\hi{गौरव}} — \emph{gaurav} — \textquotedblleft{}pride, a matter of honour\textquotedblright{}.

\begin{quote}
\hi{अगर} \hi{तुम} \hi{जीतोगे} \hi{तो} \hi{यह} \hi{पूरे} \hi{गाँव} \hi{के} \hi{लिए} \hi{गौरव} \hi{की} \hi{बात} \hi{होगी।}
\end{quote}

\begin{quote}
\emph{agar tum jītoge to yah pūre gā̃v ke liye gaurav kī bāt hogī.}
\end{quote}

\begin{quote}
If you win, it will be a matter of pride for the whole village.
\end{quote}

Masculine.

\begin{cousinweb}
\textbf{\hi{गौरव}} (\emph{gaurav}) is a tatsama, borrowed from Sanskrit \emph{gaurava}, \textquotedblleft{}weight; importance, dignity\textquotedblright{} (Monier-Williams).
\end{cousinweb}

\begin{grammarlens}[title={What you've built}]
One more word for saying what happens if.
\end{grammarlens}

\subsection*{Your turn}
\begin{itemize}
\raggedright
  \item \emph{Say it:} \emph{gaurav}
  \item \emph{Say it:} \emph{gaurav}, once more
  \item \emph{Say it:} \emph{agar tum jītoge to yah pūre gā̃v ke liye gaurav kī bāt hogī.} — If you win, it will be a matter of pride for the whole village.
  \item \emph{Recall:} say the Hindi for to attain, then the Hindi for a promotion, then say \emph{gaurav} again
\end{itemize}

\begin{tcolorbox}[breakable,colback=teal!4,colframe=teal!35!black,title={Before you move on}]
What does \hi{गौरव} mean? (\textquotedblleft{}Pride, a matter of honour\textquotedblright{}.) Say it once more.
\end{tcolorbox}
